\textbf{M4} cuts the event sequence into consecutive windows of width $\Delta t$ and
aggregates each window into a weighted undirected snapshot whose weights are event
counts. Windows with fewer than 100 events (PS; this removes nights) or 50 events (CM) are
discarded, and each snapshot contains only the nodes active in its window. The main
settings are $\Delta t=15$\,min for PS and 7\,days for CM. Every snapshot and the
static aggregate are partitioned with Leiden \citep{traag2019louvain} maximising
modularity with resolution $\gamma$ \citep{reichardt2006statistical}. A multislice
partition \citep{mucha2010community} with interslice coupling $\omega=1$ gives labels
that are consistent across windows. It is used only for tracking, because the coupling
itself draws consecutive partitions together.

Partitions are compared by normalized mutual information (NMI) \citep{danon2005comparing}
on the nodes they share. For each window, the NMI is recorded with the previous window
(for consecutive windows only, not across a night), with the static partition and, for PS,
with the classes (teachers excluded). As a noise floor, every window is partitioned a
second time with a different random seed. All measures are recomputed on 20 P1 samples.
P1 preserves the static graph, and hence the static partition, as well as the number of
events per window; H5 therefore predicts that real snapshots agree less with each other
and with the static partition than P1 snapshots do. The analysis is repeated for $\gamma\in\{0.5,1,2\}$ and
$\Delta t\in\{5,15,60\}$\,min (PS) or $\{1,7,30\}$\,days (CM).
